% GENERATED FILE. Edit canonical lessons, then run npm run generate:books.
% canonical-source-hash: fnv1a64:89d0cbe26876be14
% canonical-lessons: ML-C370-moli, ML-C370-drksaksi, ML-C370-kallam, ML-C370-moli-nalkuka, ML-C370-satyam-ceyyuka

\chapter{Testimony}
\label{ch:ml-testimony}
\hlchaptermodality{370}

\begin{chapteropening}
\textbf{By the end of this chapter:} \emph{I can say testimony, an eyewitness, a lie, to testify and to swear an oath.}

Complete the last lesson of chapter 370: I can say testimony, an eyewitness, a lie, to testify and to swear an oath.
\end{chapteropening}

\section[\texorpdfstring{\ml{മൊഴി} (moḻi)}{മൊഴി (moḻi)}]{\ml{മൊഴി} (moḻi) — testimony; a statement}

\label{lesson:ML-C370-moli}



\begin{quote}
Before the new one: say the Malayalam for an innocent person, then the Malayalam for to pass sentence.
\end{quote}

\subsection*{You'll want to know: \ml{മൊഴി}}
\textbf{\ml{മൊഴി}} — \emph{moḻi} — \textquotedblleft{}testimony; a statement\textquotedblright{}.

\begin{quote}
\ml{പോലീസ്} \ml{സാക്ഷിയുടെ} \ml{മൊഴി} \ml{എടുത്തു}.
\end{quote}

\begin{quote}
\emph{pōlīsŭ sākṣiyuṭe moḻi eṭuttu.}
\end{quote}

\begin{quote}
The police took the witness's statement.
\end{quote}

The police take it: \ml{മൊഴി} \ml{എടുക്കുക}. In old poetry it is any word or speech.

\begin{cousinweb}
\textbf{\ml{മൊഴി}} (\emph{moḻi}) is native Dravidian: its Tamil twin \emph{moḻi}, \textquotedblleft{}to say, to speak; a word, a language\textquotedblright{}, stands in DEDR 4989. In the police station it is the word a witness gives.
\end{cousinweb}

\begin{grammarlens}[title={What you've built}]
One more word for telling what happened.
\end{grammarlens}

\subsection*{Your turn}
\begin{itemize}
\raggedright
  \item \emph{Say it:} \emph{moḻi}
  \item \emph{Say it:} \emph{moḻi}, once more
  \item \emph{Say it:} \emph{pōlīsŭ sākṣiyuṭe moḻi eṭuttu.} — The police took the witness's statement.
  \item \emph{Recall:} say the Malayalam for an innocent person, then the Malayalam for to pass sentence, then say \emph{moḻi} again
\end{itemize}

\begin{tcolorbox}[breakable,colback=teal!4,colframe=teal!35!black,title={Before you move on}]
What does \ml{മൊഴി} mean? (\textquotedblleft{}Testimony; a statement\textquotedblright{}.) Say it once more.
\end{tcolorbox}
\section[\texorpdfstring{\ml{ദൃക്സാക്ഷി} (dṛksākṣi)}{ദൃക്സാക്ഷി (dṛksākṣi)}]{\ml{ദൃക്സാക്ഷി} (dṛksākṣi) — an eyewitness}

\label{lesson:ML-C370-drksaksi}



\begin{quote}
Before the new one: say the Malayalam for to pass sentence, then the Malayalam for testimony.
\end{quote}

\subsection*{You'll want to know: \ml{ദൃക്സാക്ഷി}}
\textbf{\ml{ദൃക്സാക്ഷി}} — \emph{dṛksākṣi} — \textquotedblleft{}an eyewitness\textquotedblright{}.

\begin{quote}
\ml{അപകടത്തിന്} \ml{രണ്ട്} \ml{ദൃക്സാക്ഷികൾ} \ml{ഉണ്ടായിരുന്നു}.
\end{quote}

\begin{quote}
\emph{apakaṭattinŭ raṇṭŭ dṛksākṣikaḷ uṇṭāyirunnu.}
\end{quote}

\begin{quote}
There were two eyewitnesses to the accident.
\end{quote}

Plural: \ml{ദൃക്സാക്ഷികൾ}, \emph{dṛksākṣikaḷ}.

\ml{ദൃക്}, seeing, + \ml{സാക്ഷി}, a witness: one who saw it happen.

\begin{cousinweb}
\textbf{\ml{ദൃക്സാക്ഷി}} (\emph{dṛksākṣi}) is built from Sanskrit \emph{dṛk}, the form of \emph{dṛś}, \textquotedblleft{}sight, seeing\textquotedblright{}, used in compounds (Monier-Williams), and \ml{സാക്ഷി} (\emph{sākṣi}), \textquotedblleft{}a witness\textquotedblright{}, from Sanskrit \emph{sākṣin} (CDIAL 13321).
\end{cousinweb}

\begin{grammarlens}[title={What you've built}]
One more word for telling what happened.
\end{grammarlens}

\subsection*{Your turn}
\begin{itemize}
\raggedright
  \item \emph{Say it:} \emph{dṛksākṣi}
  \item \emph{Say it:} \emph{dṛksākṣi}, once more
  \item \emph{Say it:} \emph{apakaṭattinŭ raṇṭŭ dṛksākṣikaḷ uṇṭāyirunnu.} — There were two eyewitnesses to the accident.
  \item \emph{Recall:} say the Malayalam for to pass sentence, then the Malayalam for testimony, then say \emph{dṛksākṣi} again
\end{itemize}

\begin{tcolorbox}[breakable,colback=teal!4,colframe=teal!35!black,title={Before you move on}]
What does \ml{ദൃക്സാക്ഷി} mean? (\textquotedblleft{}An eyewitness\textquotedblright{}.) Say it once more.
\end{tcolorbox}
\section[\texorpdfstring{\ml{കള്ളം} (kaḷḷaṁ)}{കള്ളം (kaḷḷaṁ)}]{\ml{കള്ളം} (kaḷḷaṁ) — a lie; a falsehood}

\label{lesson:ML-C370-kallam}



\begin{quote}
Before the new one: say the Malayalam for testimony, then the Malayalam for an eyewitness.
\end{quote}

\subsection*{You'll want to know: \ml{കള്ളം}}
\textbf{\ml{കള്ളം}} — \emph{kaḷḷaṁ} — \textquotedblleft{}a lie; a falsehood\textquotedblright{}.

\begin{quote}
\ml{സാക്ഷി} \ml{കോടതിയിൽ} \ml{കള്ളം} \ml{പറഞ്ഞു}.
\end{quote}

\begin{quote}
\emph{sākṣi kōṭatiyil kaḷḷaṁ paṟaññu.}
\end{quote}

\begin{quote}
The witness told a lie in court.
\end{quote}

To lie is \ml{കള്ളം} \ml{പറയുക}, the same as \ml{നുണ} \ml{പറയുക}. \ml{കള്ളൻ}, a thief, comes from the same root.

\begin{cousinweb}
\textbf{\ml{കള്ളം}} (\emph{kaḷḷaṁ}) is native Dravidian: DEDR 1372 lists Malayalam \emph{kaḷḷam}, \textquotedblleft{}theft, untruth, false\textquotedblright{}, from the root that also gives \ml{കള്ളൻ} (\emph{kaḷḷan}), \textquotedblleft{}a thief\textquotedblright{}.
\end{cousinweb}

\begin{grammarlens}[title={What you've built}]
One more word for telling what happened.
\end{grammarlens}

\subsection*{Your turn}
\begin{itemize}
\raggedright
  \item \emph{Say it:} \emph{kaḷḷaṁ}
  \item \emph{Say it:} \emph{kaḷḷaṁ}, once more
  \item \emph{Say it:} \emph{sākṣi kōṭatiyil kaḷḷaṁ paṟaññu.} — The witness told a lie in court.
  \item \emph{Recall:} say the Malayalam for testimony, then the Malayalam for an eyewitness, then say \emph{kaḷḷaṁ} again
\end{itemize}

\begin{tcolorbox}[breakable,colback=teal!4,colframe=teal!35!black,title={Before you move on}]
What does \ml{കള്ളം} mean? (\textquotedblleft{}A lie; a falsehood\textquotedblright{}.) Say it once more.
\end{tcolorbox}
\section[\texorpdfstring{\ml{മൊഴി} \ml{നൽകുക} (moḻi nalkuka)}{മൊഴി നൽകുക (moḻi nalkuka)}]{\ml{മൊഴി} \ml{നൽകുക} (moḻi nalkuka) — to testify; to give a statement}

\label{lesson:ML-C370-moli-nalkuka}



\begin{quote}
Before the new one: say the Malayalam for an eyewitness, then the Malayalam for a lie.
\end{quote}

\subsection*{You'll want to know: \ml{മൊഴി} \ml{നൽകുക}}
\textbf{\ml{മൊഴി} \ml{നൽകുക}} — \emph{moḻi nalkuka} — \textquotedblleft{}to testify; to give a statement\textquotedblright{}.

\begin{quote}
\ml{ദൃക്സാക്ഷി} \ml{പോലീസിന്} \ml{മൊഴി} \ml{നൽകി}.
\end{quote}

\begin{quote}
\emph{dṛksākṣi pōlīsinŭ moḻi nalki.}
\end{quote}

\begin{quote}
The eyewitness gave a statement to the police.
\end{quote}

Past: \ml{മൊഴി} \ml{നൽകി}, \emph{moḻi nalki}.

Who receives it takes -\ml{ക്ക്} or -\ml{ന്}: \ml{പോലീസിന്} \ml{മൊഴി} \ml{നൽകി}.

\begin{cousinweb}
\textbf{\ml{മൊഴി} \ml{നൽകുക}} (\emph{moḻi nalkuka}) joins \ml{മൊഴി} (\emph{moḻi}), \textquotedblleft{}a statement\textquotedblright{}, a native Dravidian word (DEDR 4989), and \ml{നൽകുക}, \textquotedblleft{}to give\textquotedblright{}.
\end{cousinweb}

\begin{grammarlens}[title={What you've built}]
One more word for telling what happened.
\end{grammarlens}

\subsection*{Your turn}
\begin{itemize}
\raggedright
  \item \emph{Say it:} \emph{moḻi nalkuka}
  \item \emph{Say it:} \emph{moḻi nalkuka}, once more
  \item \emph{Say it:} \emph{dṛksākṣi pōlīsinŭ moḻi nalki.} — The eyewitness gave a statement to the police.
  \item \emph{Recall:} say the Malayalam for an eyewitness, then the Malayalam for a lie, then say \emph{moḻi nalkuka} again
\end{itemize}

\begin{tcolorbox}[breakable,colback=teal!4,colframe=teal!35!black,title={Before you move on}]
What does \ml{മൊഴി} \ml{നൽകുക} mean? (\textquotedblleft{}To testify; to give a statement\textquotedblright{}.) Say it once more.
\end{tcolorbox}
\section[\texorpdfstring{\ml{സത്യം} \ml{ചെയ്യുക} (satyaṁ ceyyuka)}{സത്യം ചെയ്യുക (satyaṁ ceyyuka)}]{\ml{സത്യം} \ml{ചെയ്യുക} (satyaṁ ceyyuka) — to swear an oath}

\label{lesson:ML-C370-satyam-ceyyuka}



\begin{quote}
Before the new one: say the Malayalam for a lie, then the Malayalam for to testify.
\end{quote}

\subsection*{You'll want to know: \ml{സത്യം} \ml{ചെയ്യുക}}
\textbf{\ml{സത്യം} \ml{ചെയ്യുക}} — \emph{satyaṁ ceyyuka} — \textquotedblleft{}to swear an oath\textquotedblright{}.

\begin{quote}
\ml{സാക്ഷി} \ml{കോടതിയിൽ} \ml{സത്യം} \ml{ചെയ്തു}.
\end{quote}

\begin{quote}
\emph{sākṣi kōṭatiyil satyaṁ ceytu.}
\end{quote}

\begin{quote}
The witness swore an oath in court.
\end{quote}

Past: \ml{സത്യം} \ml{ചെയ്തു}, \emph{satyaṁ ceytu}.

More than telling the truth: promising by oath that you are. A new verb, though both parts are known.

\begin{cousinweb}
\textbf{\ml{സത്യം} \ml{ചെയ്യുക}} (\emph{satyaṁ ceyyuka}) joins \ml{സത്യം} (\emph{satyaṁ}), borrowed from Sanskrit \emph{satya}, \textquotedblleft{}true\textquotedblright{}, and also \textquotedblleft{}a promise, an oath\textquotedblright{} (CDIAL 13112), and \ml{ചെയ്യുക}, \textquotedblleft{}to do\textquotedblright{}: to make an oath.
\end{cousinweb}

\begin{grammarlens}[title={What you've built}]
One more word for telling what happened.
\end{grammarlens}

\subsection*{Your turn}
\begin{itemize}
\raggedright
  \item \emph{Say it:} \emph{satyaṁ ceyyuka}
  \item \emph{Say it:} \emph{satyaṁ ceyyuka}, once more
  \item \emph{Say it:} \emph{sākṣi kōṭatiyil satyaṁ ceytu.} — The witness swore an oath in court.
  \item \emph{Recall:} say the Malayalam for a lie, then the Malayalam for to testify, then say \emph{satyaṁ ceyyuka} again
\end{itemize}

\begin{tcolorbox}[breakable,colback=teal!4,colframe=teal!35!black,title={Before you move on}]
What does \ml{സത്യം} \ml{ചെയ്യുക} mean? (\textquotedblleft{}To swear an oath\textquotedblright{}.) Say it once more.
\end{tcolorbox}
